This chapter provides an in-depth analysis and interpretation of the experimental results presented in Chapter 4. The discussion contextualizes the findings within the broader field, relates them back to the research objectives, and compares the performance of the proposed framework against existing approaches.

\section{Findings Analysis}
\label{sec:findings_analysis}

The experimental results demonstrate the effectiveness of the proposed adaptive framework integrating memory forensics, computer vision, and machine learning for detecting fileless and polymorphic malware. Several key findings warrant further analysis:

\begin{itemize}
    \item \textbf{Superior Performance of Linear SVM:} In the multi-class classification task (Phase 1), Linear SVM achieved the highest accuracy (97.49\%) among the tested classifiers (Table \ref{tab:classification_comparison}). While other complex models like SMO (RBF) also performed well, the success of Linear SVM suggests that the high-dimensional feature space created by fusing GIST and HOG features might be linearly separable to a large extent, or at least that a linear boundary is sufficient to achieve excellent separation for most classes. Linear SVMs are often effective in high-dimensional spaces, especially when the number of dimensions is comparable to or larger than the number of samples, and they can be less prone to overfitting than more complex models like kernel SVMs or deep networks in such scenarios, particularly with appropriate regularization (C=10).

    \item \textbf{Effectiveness of GIST+HOG Feature Fusion:} The feature ablation study (Table \ref{tab:ablation_study}) provided strong quantitative evidence supporting the hypothesis that combining global (GIST) and local (HOG) visual features yields a more powerful representation than using either alone. The significant average accuracy improvement (+3.16\%) achieved with the fused features underscores the complementary nature of these descriptors. GIST captures the overall layout and texture, while HOG focuses on finer-grained edge and shape information derived from byte patterns. Their combination allows the classifier to leverage both macro- and micro-level visual cues present in the memory dump images, leading to more robust and accurate classification. This validates the feature engineering approach as a key contributor to the framework's success.

    \item \textbf{Importance of Image Representation (Column Width):} The column width analysis (Table \ref{tab:col_width_impact}) revealed that the way the 1D memory byte stream is mapped to a 2D image significantly impacts classification performance. The superior results obtained with the 4096px width suggest that this configuration better preserves inherent structural patterns within the memory data compared to narrower or square-root based widths. This finding implies that spatial relationships between bytes, as captured in the 2D visualization, contain discriminative information, and the choice of image width is critical for making these patterns accessible to the feature extraction algorithms. It supports the core premise of the vision-based approach – that visual structure in binary data is meaningful for malware analysis.

    \item \textbf{Transformative Impact of UMAP on Unknown Malware Detection:} Perhaps the most striking finding is the dramatic improvement in detecting unknown malware families after applying UMAP (Tables \ref{tab:umap_improvement} and \ref{tab:fold_umap_analysis}). The accuracy leap of up to 36.98\% in challenging folds demonstrates UMAP's ability to create a lower-dimensional embedding (4-dimensions) where novel malware samples are more effectively distinguished from benign ones. Critically, UMAP not only boosted overall accuracy but also drastically improved the precision for benign samples (from as low as 0.24 to 0.71), significantly reducing the false positive rate which was a major limitation pre-UMAP. This suggests that supervised UMAP, by leveraging class labels during reduction, successfully learned a transformation that pushes benign and malicious clusters further apart in the embedding space, even accommodating the structure of unseen malware families within the malicious cluster boundary. The preservation of both local and global data structure by UMAP \citep{mcinnes2018umap} appears crucial for this success, allowing the model to generalize better than when operating in the original high-dimensional space.

    \item \textbf{Necessity of Hyperparameter Tuning:} The UMAP configuration analysis (Table \ref{tab:umap_config_analysis}) underscores the importance of careful hyperparameter optimization. The optimal performance achieved with `n_neighbors=55`, `min_dist=1.0`, and Manhattan distance highlights that UMAP's effectiveness is sensitive to these settings, which control how the algorithm balances local versus global structure preservation and the density of the resulting embedding. Finding the right parameters through systematic experimentation was key to unlocking the full potential of UMAP for this specific task and dataset.
\end{itemize}

Collectively, these findings validate the core components of the proposed methodology. The combination of RGB memory visualization, fused GIST+HOG features, appropriate classifier selection (Linear SVM), and optimized UMAP-based dimensionality reduction provides a robust and effective solution for detecting known malware families with high accuracy and significantly enhancing the challenging task of identifying previously unseen threats.

\section{Comparison with Existing Approaches}
\label{sec:comparison_literature}

The performance achieved by the proposed framework can be contextualized by comparing it with relevant existing work, particularly the baseline study by \citet{bozkir2021catch} which employed a similar methodology involving memory visualization and GIST/HOG features.

\begin{itemize}
    \item \textbf{Multi-Class Classification Accuracy:} The proposed framework achieved a peak multi-class accuracy of 97.49\% using Linear SVM (Table \ref{tab:classification_comparison}). This significantly surpasses the peak accuracy of 96.39\% (using SMO) reported by \citet{bozkir2021catch}, and specifically improves upon their Linear SVM result (93.26\%). This suggests that the use of RGB images (compared to potential grayscale in some interpretations of related work) combined with the specific GIST+HOG feature extraction parameters and potentially the more contemporary dataset used in this thesis, contributes to a more effective classification model. While other classifiers like SMO showed comparable performance to the baseline, the advantage demonstrated with Linear SVM is notable.

    \item \textbf{Feature Fusion Benefit:} The ablation study (Table \ref{tab:ablation_study}) showed an average accuracy gain of +3.16\% from fusing GIST and HOG features compared to using GIST alone. This reinforces the findings of \citet{bozkir2021catch} who also observed benefits from fusion, although the improvement margin observed in this thesis (+3.16\%) appears slightly larger than the ~2.3\% reported in their work (based on comparing their combined vs GIST-only results implicitly), suggesting the fusion strategy or feature extraction details here might be more effective.

    \item \textbf{Unknown Malware Detection Enhancement via UMAP:} The application of UMAP resulted in accuracy improvements of up to +36.98\% (Table \ref{tab:fold_umap_analysis}) in detecting unknown malware families. This significantly exceeds the maximum improvement of +21.83\% reported by \citet{bozkir2021catch} from their initial application of UMAP for the same purpose. This substantial difference highlights the effectiveness of the specific supervised UMAP approach, the optimized hyperparameters (`n_neighbors=55`, `min_dist=1.0`, Manhattan distance), and potentially the nature of the dataset used in this work. It strongly supports the conclusion that well-tuned UMAP is a highly potent technique for enhancing zero-day malware detection in memory-visualization-based frameworks.

    \item \textbf{Alignment with General Trends:} The overall success aligns with the broader trends identified in the literature review (Chapter 2). It confirms the value of memory forensics for detecting evasive threats \citep{kara2022fileless, nissim2019volatile}, supports the viability of vision-based approaches initiated by works like \citet{nataraj2011malware}, and demonstrates the crucial role of machine learning \citep{gibert2020rise, khalid2023insight} and advanced dimensionality reduction techniques \citep{sharma2018efficacy} in building effective detection systems. The high performance achieved on a dataset containing contemporary fileless and polymorphic families (like Kovter, Poweliks, Emotet, Trickbot) further strengthens the practical relevance compared to studies potentially using older datasets.
\end{itemize}

In summary, the proposed framework not only achieves high accuracy comparable to or exceeding state-of-the-art methods but also demonstrates significant advancements, particularly in the fusion of GIST+HOG features and the highly effective application of optimized UMAP for enhancing the detection of unknown malware. These comparisons validate the contributions of this thesis in advancing memory-based malware detection capabilities.
